% 00 — Uvod.

\section{Uvod}
\label{sec:uvod}

Ovaj seminarski rad obrađuje četiri zadatka iz kolegija \emph{Optimizacijske
metode u razvoju i upravljanju sustava} \cite{jokic_zadatak}. Iako se zadaci na
prvi pogled bave različitim temama --- od kvadratnog programa s dva ograničenja
do stabilizacije hibridnog dinamičkog sustava --- povezuje ih zajednička nit:
svaki se svodi na \emph{konveksan} optimizacijski problem, za koji lokalni
optimum ujedno je i globalni, a uvjeti optimalnosti su i nužni i dovoljni
\cite[str.~61]{jokic_konveksnost}.

Poglavlje~\ref{sec:zad1} obrađuje kvadratni program s linearnim ograničenjima.
Analizira se konveksnost, postavljaju se i analitički rješavaju
Karush--Kuhn--Tuckerovi uvjeti \cite[str.~59]{jokic_kkt}, rezultat se provjerava
funkcijom \texttt{quadprog} i modeliranjem u YALMIP-u, te se postavlja i rješava
Lagrangeov dualni problem. Na kraju se pokazuje kako Lagrangeovi množitelji
mjere osjetljivost optimalne vrijednosti na perturbacije ograničenja.

Poglavlje~\ref{sec:zad2} isti problem promatra uz nesigurne koeficijente
ograničenja. Pokazuje se da se zahtjev dopustivosti za sve realizacije svodi na
konačan skup ograničenja najgoreg slučaja i da robusna protuformulacija ostaje
konveksna, pa se rješava istim alatom. Kvantificira se i cijena robusnosti.

Poglavlje~\ref{sec:zad3} bavi se oblikovanjem profila ceste na neravnom terenu.
Naglasak je na \emph{formulaciji}: pokazuje se da izbor mjere troška određuje
klasu problema, te da diskretizacija derivacija podijeljenim razlikama zadržava
linearnost i vodi na linearni program.

Poglavlje~\ref{sec:zad4} prelazi na dinamičke sustave. Analizira se prekidački
sustav s dva stabilna moda koji je unatoč tome nestabilan, i sintetizira se
statički regulator koji ga stabilizira. Sinteza se oslanja na zajedničku
kvadratnu Ljapunovljevu funkciju \cite[str.~4]{jokic_ljapunov}, čime se problem
upravljanja svodi na sustav linearnih matričnih nejednakosti.

Svi numerički rezultati dobiveni su u MATLAB-u, po jednom skriptom za svaki
zadatak. Skripte su u cijelosti navedene u prilogu (poglavlje~\ref{sec:kod}), a
nalaze se i u direktoriju \texttt{src/} repozitorija; reproduciraju sve brojeve
i slike iz ovog rada.
